\documentclass[9pt]{extarticle}
\usepackage[margin=1in, top=0.85in, bottom=0.85in]{geometry}
\usepackage{hyperref}
\usepackage{booktabs}
\usepackage{enumitem}
\usepackage{titlesec}
\usepackage{multicol}
\usepackage{mdframed}

\setlength{\parskip}{3pt}
\setlength{\parindent}{0pt}
\setlength{\multicolsep}{4pt}
\setlength{\columnsep}{14pt}

\titlespacing*{\section}{0pt}{10pt}{5pt}
\titlespacing*{\subsection}{0pt}{8pt}{3pt}

\setlist{noitemsep, topsep=2pt, partopsep=0pt, parsep=0pt}

\begin{document}

\begin{center}
    {\large\textbf{CS 581 --- Nuclear Cybersecurity}} \\[2pt]
    {\normalsize\textbf{Companion Document: AI Coding Assistants in This Course}} \\[1pt]
    {\normalsize Boise State University \quad|\quad Fall 2026}
\end{center}
\vspace{4pt}

%------------------------------------------------------------
\section*{Why This Course Requires an AI Coding Assistant}
%------------------------------------------------------------

This course requires every student to use an AI Coding Assistant throughout the
semester. This is not a convenience --- it is a deliberate professional choice.

Nuclear cybersecurity practitioners at national laboratories, utilities, and
regulatory agencies are already using AI-assisted tools to draft threat models,
analyze regulatory documents, triage incidents, and produce security assessments
at greater depth than was previously possible. A graduate course that ignores
this reality would be preparing you for a profession that no longer exists.

More importantly: the capabilities covered in this document --- thinking modes,
function calling, structured output, multimodal analysis, and MCP tool
integration --- are not features. \textbf{They are the skills this course builds.}
Each workshop is designed around one or more of these capabilities applied to a
real nuclear security problem.

%------------------------------------------------------------
\section*{Choosing Your Tool}
%------------------------------------------------------------

Selecting your AI Coding Assistant is the first real decision of this course ---
and it is itself an assignment. You will evaluate and justify your choice based
on the capability dimensions below. You may use one tool for the entire semester
or combine tools by task; either approach is valid if you can defend it.

\textbf{Evaluate any candidate tool against these capability areas:}

\begin{itemize}
    \item \textbf{Thinking \& Reasoning Depth} --- Does it support configurable
          reasoning modes (deep vs.\ fast)? Can it reason between tool calls?
          Does it preserve reasoning across turns?
    \item \textbf{Context Length} --- How much can it hold in a single session?
          Loading full regulatory documents requires significant context capacity.
    \item \textbf{Function Calling \& Tool Use} --- Can it invoke external tools,
          APIs, or scripts during a session? This is essential for live data workflows.
    \item \textbf{Structured Output} --- Can it reliably return JSON or other
          defined formats? Required for compliance matrices and threat models.
    \item \textbf{Vision \& Multimodal Input} --- Can it analyze images, diagrams,
          and PDFs? Critical for architecture review and document analysis.
    \item \textbf{MCP / Tool Integration} --- Does it support Model Context Protocol
          servers for web search, document reading, and live data access?
    \item \textbf{Provenance \& Trust} --- Who built it, where is data processed,
          and is it appropriate for use with sensitive security content?
\end{itemize}

The last criterion is non-negotiable in nuclear security contexts. Only tools
developed and hosted by US-based providers with clear data handling policies
are appropriate for this course.

Regardless of tool, you are responsible for the accuracy of everything you submit.
Verify all regulatory citations, technical claims, and system-specific assertions
against primary sources.

%------------------------------------------------------------
\section*{Core Capabilities --- The Skills You Are Building}
%------------------------------------------------------------

The following capabilities are not optional features to explore. They are the
technical foundation of AI-assisted security work. You will use each of them
in the course workshops.

\subsection*{1. Thinking Modes \& Deep Reasoning}

Modern AI coding assistants support configurable reasoning depth. Understanding
when and how to engage deep thinking is itself a professional skill.

\medskip

\textbf{What it is:} Deep Thinking enables Chain-of-Thought reasoning --- the model
works through a problem systematically before producing an answer. This
significantly improves accuracy on complex, multi-step analysis tasks.

\medskip

\textbf{Modes to look for in your chosen tool:}
\begin{itemize}
    \item \texttt{enabled} --- Model auto-decides when to think deeply (default)
    \item \texttt{disabled} --- Direct answers, lower latency, for simple queries
    \item \textbf{Interleaved Thinking} --- Reasoning occurs between tool calls,
          allowing the model to interpret each result before deciding the next step
    \item \textbf{Preserved Thinking} --- Retains reasoning from prior turns;
          default on Coding Plan endpoints; critical for multi-session workflows
    \item \texttt{reasoning\_effort} --- Scales from \texttt{none} to \texttt{max};
          use \texttt{max} for threat modeling, regulatory analysis, and complex
          incident scenarios
\end{itemize}

\textbf{Nuclear security application:} Enable deep thinking (\texttt{max} effort)
for threat modeling, attack chain analysis, and regulatory gap assessment. Disable
for simple lookups, format conversions, or quick summaries. Misjudging this
wastes tokens and time --- or produces shallow analysis on complex problems.

\medskip

\textbf{Critical implementation note:} When using preserved or interleaved
thinking, always return the complete, unmodified \texttt{reasoning\_content}
back to the model. Truncating it degrades performance and breaks cache efficiency.

\subsection*{2. Function Calling \& Tool Use}

Function calling allows the AI to interact with external systems --- databases,
APIs, file systems, and security tools --- during a session. This transforms
the model from a text generator into an active analysis agent.

\medskip

\textbf{What it is:} You define callable functions (tools) with names, descriptions,
and parameter schemas. The model decides when to invoke them, passes structured
arguments, receives results, and reasons over the output before continuing.

\medskip

\textbf{Security-relevant applications:}
\begin{itemize}
    \item Query a CVE database during a vulnerability analysis session
    \item Execute a network scan script and have the model interpret results
    \item Pull NRC regulatory text from an API and map it to a control architecture
    \item Retrieve ICS/SCADA log data and run anomaly detection routines
    \item Chain multiple operations: search $\rightarrow$ analyze $\rightarrow$
          draft recommendation
\end{itemize}

\textbf{Best practices:} Single responsibility per function. Validate all inputs.
Log all calls. Implement permission controls. Never pass unvalidated user input
directly to a function that executes code or queries a database.

\textbf{Nuclear security application:} In the Monitoring \& Anomaly Detection
workshop, you will build a function-calling workflow that queries a simulated
ICS data stream, flags anomalies, and drafts an incident notification. This
is the same pattern used in real SOC environments.

\subsection*{3. Structured Output}

Structured output forces the model to return data in a defined format ---
typically JSON --- rather than free-form prose. This is essential when your
output needs to be machine-readable, consistently formatted, or fed into
another system.

\medskip

\textbf{What it is:} Set \texttt{response\_format: \{"type": "json\_object"\}}
and define your expected schema in the system prompt. The model guarantees
output conformance, enabling reliable programmatic processing.

\medskip

\textbf{Security-relevant applications:}
\begin{itemize}
    \item Threat models as structured JSON (actor, vector, likelihood, impact, mitigation)
    \item Compliance matrices with pass/fail fields per regulatory requirement
    \item Incident timelines with typed event fields
    \item Vulnerability assessments with CVSS-style scoring fields
    \item Configuration audit results with remediation status flags
\end{itemize}

\textbf{Best practices:} Start with simple schemas. Name fields explicitly.
Add descriptions for each field in the system prompt. Implement validation
against the schema after receiving output --- JSON mode reduces but does not
eliminate malformed responses in complex scenarios.

\textbf{Nuclear security application:} Your compliance matrix workshop (NRC 10 CFR
73.54) will require structured output. A JSON-formatted compliance matrix can be
imported directly into a spreadsheet or ticketing system --- exactly how this
work is done at a licensed facility.

\subsection*{4. Vision \& Multimodal Analysis}

Vision Language Models (VLMs) accept images, diagrams, documents, and in some
cases audio and video as inputs alongside text. This capability is particularly
powerful for security work where evidence is rarely just text.

\medskip

\textbf{Modalities to look for in your chosen tool:}
\begin{itemize}
    \item \textbf{Images} --- Architecture diagrams, network maps, screenshots,
          scanned documents
    \item \textbf{Files / PDFs} --- Regulatory documents, configuration exports,
          inspection reports
    \item \textbf{Audio} --- Transcription of interviews or practitioner briefings
    \item \textbf{OCR} --- Structured text extraction from scanned documents
\end{itemize}

\textbf{Security-relevant applications:}
\begin{itemize}
    \item Upload a network architecture diagram and ask the model to identify
          attack surfaces and missing segmentation controls
    \item Feed a scanned ICS schematic and map it to the Purdue Model zones
    \item Analyze a screenshot of anomalous SCADA readings against baseline
    \item Extract structured data from a scanned NRC inspection report
\end{itemize}

\textbf{Nuclear security application:} In the Physical Protection workshop, you will
upload a facility diagram and use a VLM to identify cyber-physical intersection
points --- where a digital compromise could affect a physical safety boundary.

\subsection*{5. Model Context Protocol (MCP) Integration}

MCPs extend the model's reach beyond its training data by connecting it to live
external tools and data sources during a session.

\medskip

\textbf{Common MCP servers to look for:}
\begin{itemize}
    \item \textbf{Web Search} --- Real-time queries for current threat intelligence,
          recent CVEs, and regulatory updates
    \item \textbf{Web Reader} --- Full-page content extraction from URLs; feed the
          model a live NRC advisory or CISA alert directly
    \item \textbf{Vision} --- On-demand image analysis within a session
    \item \textbf{Document Reader} --- Structured extraction from PDFs and reports
\end{itemize}

\textbf{Security-relevant applications:}
\begin{itemize}
    \item Pull the latest CISA ICS advisory and have the model assess applicability
          to a specific nuclear plant architecture
    \item Search for recent threat actor TTPs and incorporate them into a live
          threat model
    \item Read an NRC regulatory document and extract all applicable control
          requirements automatically
\end{itemize}

\textbf{Nuclear security application:} The Insider Threat workshop will use Web
Search + Web Reader MCPs to pull current case studies and have the model
synthesize a threat profile grounded in recent real-world incidents --- not
just training data.

%------------------------------------------------------------
\section*{Context Length \& Session Management}
%------------------------------------------------------------

Leading AI coding assistants support up to \textbf{1 million tokens of context}. This is not a
trivial capability --- it means you can load entire regulatory frameworks,
multi-document collections, or long session histories into a single working
context. Used well, it eliminates the need to re-explain your problem across
sessions.

\medskip

\textbf{Practical guidance:}
\begin{itemize}
    \item Load full regulatory documents at session start rather than pasting excerpts
    \item Keep your session log as a running context --- preserved thinking benefits
          from continuity across a multi-step analysis
    \item Be aware that very large contexts increase latency and token cost; use
          context caching where available
    \item Maximum output tokens vary by model tier --- plan your prompts accordingly
          for long deliverables like compliance matrices or threat models
\end{itemize}

%------------------------------------------------------------
\section*{Academic Integrity in an AI-Assisted Course}
%------------------------------------------------------------

Using an AI Coding Assistant is \textbf{required}, not prohibited. However, the
following remain your sole responsibility:

\begin{itemize}
    \item The accuracy of all regulatory citations and technical claims
    \item The analysis, judgment, and conclusions in your work products
    \item Understanding your submission well enough to discuss and defend it
\end{itemize}

Session logs are submitted with every assignment. They make your reasoning
process visible. If you cannot explain or defend your submitted work in a
follow-up discussion, it will not be counted.

%------------------------------------------------------------
\section*{Building Your Portfolio on GitHub}
%------------------------------------------------------------

Every workshop produces a real security artifact. Over the semester you will
build: threat models, access control policies, protocol analyses, compliance
matrices, monitoring plans, an incident response package, and a capstone project.

\medskip

\textbf{Repository structure:}
\begin{itemize}
    \item \texttt{/workshops} --- one folder per module; artifact + session log
    \item \texttt{/project} --- semester project materials and interim drafts
    \item \texttt{/written} --- ethics paper and final written analysis
    \item \texttt{README.md} --- professional summary of the portfolio
\end{itemize}

A well-maintained repository is something you can link from a resume or share
directly with a hiring manager. The combination of security artifacts and
AI-assisted workflows is a differentiator that very few candidates currently have.

\vspace{8pt}
\hrule
\vspace{4pt}
{\small \textit{CS 581 --- Nuclear Cybersecurity \quad|\quad
Christopher Spirito \quad|\quad
\href{mailto:cspirito@boisestate.edu}{cspirito@boisestate.edu} \quad|\quad
Fall 2026}}

\end{document}
